\chapter{Digest notes}
\label{chap:notes}

A running list of questions about the manuscript, collected while writing
the blueprint. Each item points to the node where it arises. Items are removed
once resolved; record the resolution in the relevant node.

\section*{Gaps and unverified steps}
\begin{enumerate}
  \item \textbf{Mixed pairing characters}
        (Proposition~\ref{prop:mixed-paired-flux}). The Melnikov and
        pairing-limit statements cover only $\pm1$, not the $\pm i$
        characters needed here. The local endpoint formulas used for the
        weighted transport gains already cover all four characters, but do
        not by themselves extend the canonical pairing limit.
  \item \textbf{Period asymptotics} (Lemma~\ref{lem:period-asymptotics}).
        The sharp value and derivative remainders still need a precise
        elliptic-integral reference or a quantified proof. Differentiating
        a value $O$-term is not a proof of the derivative bound.
        The coarse value bounds are separate in Lemma~\ref{lem:period-coarse},
        from the independently proved saddle-distance estimate.
        Lemma~\ref{lem:period-derivative-bounds} gives the coarse derivative
        estimates needed for transport gains directly from the regular
        integral. The odd--odd finite parts use only the leading value
        asymptotic of Lemma~\ref{lem:period-leading-asymptotic}, also derived
        from the defining integral. The elliptic identity and sharp
        remainders are no longer sector prerequisites.
  \item \textbf{Final odd--odd interfaces}
        (Lemmas~\ref{lem:dd-bulk-picone} and~\ref{lem:dd-picone-current}).
        Bulk coercivity requires explicit finite-energy hypotheses; without
        them a divergent Bochner integral would be assigned value zero in
        Lean. The current's full-level remainder estimate is separate from
        the coefficient sign certificate.
        Only boundary nonnegativity is consumed by exclusion; strict
        positivity is unnecessary.
  \item \textbf{Certificate verification}
        (Chapter~\ref{chap:certificates}). A Bernstein proof would require
        exact polynomialization linked to the defining operators and a sound
        coefficient checker. The original verification scripts are missing;
        another proof method may be preferable.
\end{enumerate}

\section*{Possible simplifications}
\begin{enumerate}
  \item \textbf{The $B$-weighted level trace}
        (Lemma~\ref{lem:dd-boundaryB-trace-limit}). The bound
        $|B(r)|\le6\sqrt r$ cancels the quarter graph's singular time Jacobian,
        leaving a uniformly bounded density on a fixed finite interval.
        Dominated convergence replaces orbit comparison and separate corner
        tails, using only bounded continuous values. This route is proved in Lean.
  \item \textbf{Virial identity.} For smooth eigenfunctions the virial identity
        is a three-line computation (Lemma~\ref{lem:virial-smooth}). Stating
        it this way in the paper, followed by the justification on the $L^2$
        domain, would make the mechanism more visible.
  \item \textbf{Constrained weighted mean.} Completing a square gives the
        $81/80$ estimate directly from $\int wz=0$ and $1\le w\le5/4$
        (Lemma~\ref{lem:dd-constrained-poincare}). This is proved in Lean
        and removes the separate Kantorovich bound from the bulk route.
  \item \textbf{Weighted Poincar\'e without Sobolev density.}
        Cutoff Chebyshev-mode fluxes and Bessel give the estimate on the maximal
        weak domain (Lemmas~\ref{lem:zonal-weak-pairing} and~\ref{lem:zonal-poincare}).
        The weighted flux vanishes cubically at each endpoint, so only weighted
        $L^2$ is needed to remove the cutoffs. No sphere Sobolev identification
        or equality expansion of the gradient energy is required. For the weak
        pairing, a uniform bound on $|\cos a\,\chi_j'(a)|$ and pointwise
        stabilization of the cutoffs suffice for dominated convergence against
        weighted $L^1$; no quantitative endpoint-layer estimate is needed.
  \item \textbf{Picone quotient domain.} The ground state and its reciprocal
        are $C^1$ even at the center (Lemma~\ref{lem:dd-ground-c1}). Bounded
        multiplication proves the weighted quotient domain without Picone
        energy, matrix coercivity or point removal
        (Lemma~\ref{lem:dd-bulk-form-domain}); this is checked in Lean.
  \item \textbf{Second mixed sector.} It is obtained from the first by the
        rotation $R$ (Lemma~\ref{lem:sector-swap},
        Corollary~\ref{cor:mixed-other-exclusion}).
  \item \textbf{Resonant vanishing} (Lemma~\ref{lem:resonant-vanishing}). The argument
        needs neither $S'<0$ nor square integrability of $v_k$: differentiability of
        $S$ at the resonant level and $\int|v_k|\,dh<\infty$ suffice, since
        $1/|h-h_k|$ is not integrable. Similarly the good windows
        (Lemma~\ref{lem:good-windows}) need only continuity of $S$ and $S\to\infty$
        (Lemma~\ref{lem:period-basic}), not monotonicity or the coarse asymptotics.
  \item \textbf{Melnikov condition} (Lemma~\ref{lem:melnikov}). The rotation character of
        $v$ is not used, only that of $f$; the H\"older estimate is needed only on a
        neighbourhood of the closed separatrix segment, which projects injectively to the
        torus; and the boundary term needs only continuity of $f$ at the saddles.
\end{enumerate}

\section*{Presentation}
\begin{enumerate}
  \item \textbf{Notation clashes}: $r$, $s$, $h$, $w$, $G$, $V$, $D$ each have
        several meanings across sections (remarks in
        Chapters~\ref{chap:geometry} and~\ref{chap:mixed}).
\end{enumerate}

\section*{Checks done so far}
Scratch computations (not part of the repository) confirmed numerically or
symbolically:
the odd--odd and mixed field identities
(Lemmas~\ref{lem:dd-virial-algebra} and \ref{lem:mixed-field-algebra}, the
last apart from the formula for $M$);
the sums of squares and the affine splitting
(Lemmas~\ref{lem:trace-sos-smooth}, \ref{lem:affine-splitting},
\ref{lem:mixed-boundary-square}) on test polynomials;
the determinant formula of Lemma~\ref{lem:mixed-G-positive};
and the coefficient table of Lemma~\ref{cert:dd-residual} against a direct
computation of $Pw_*+\tfrac52s^2w_*$.
These are sanity checks, not proofs. The dipole identities
(Lemma~\ref{lem:dipole-identities}) and the normal field identities
(Lemma~\ref{lem:normal-field-algebra}) are now proved in Lean.
